\maketitle
\noindent\textbf{Manuscrito de trabalho: análise documental piloto assistida por IA, com revisão humana pendente.}

\begin{center}\textbf{From a student GDD to an AI-implemented game: design traceability and ambiguities in Night Heist}\end{center}
\section*{\centering Resumo}
\begin{resumotexto}
Este trabalho examina a correspondência entre um Game Design Document (GDD) elaborado por discentes e a implementação do jogo Night Heist: Escape Phase, atribuída a um agente ChatGPT pelo responsável pelo projeto. A questão investigada é como decisões documentadas de design aparecem no código e quais divergências e parâmetros não especificados surgem nessa tradução. Adota-se um estudo documental de caso único, com análise piloto de seis unidades selecionadas do GDD e inspeção estática de uma versão identificada do repositório. O instrumento registra localizador da especificação, critério observável, localização no código, categoria de correspondência, ambiguidade, proveniência da decisão e estado de revisão. Foram identificadas correspondências relativas aos quatro pavimentos, ao esgotamento do temporizador e ao limiar de exposição às câmeras; uma divergência literal entre retorno ao ponto exato e tolerância espacial; e insuficiência de especificação para comparar numericamente a duração da missão e o efeito do terminal de hack. As classificações são preliminares, sem execução recente do jogo, cobertura integral do GDD ou revisão humana concluída. A contribuição é uma matriz auditável que diferencia decisão explícita, concretização operacional e lacuna documental. O estudo não mede aprendizagem nem atribui decisões individuais aos estudantes ou ao agente na ausência de registros. A validade dos achados depende da conferência do documento original, da revisão das categorias e de verificações técnicas posteriores.

\end{resumotexto}
\noindent\textbf{Palavras-chave:} game design; inteligência artificial generativa; rastreabilidade; análise documental; desenvolvimento de jogos.

\noindent\begin{minipage}{\linewidth}
\section*{\centering Abstract}
\begin{resumotexto}
This study examines the correspondence between a student-authored Game Design Document (GDD) and the implementation of Night Heist: Escape Phase, attributed to a ChatGPT agent by the project coordinator. It asks how documented design decisions appear in code and which divergences and unspecified parameters arise during translation. A single-case documentary study is adopted, with a pilot analysis of six selected GDD units and static inspection of an identified repository version. The instrument records the specification locator, observable criterion, code location, correspondence category, ambiguity, decision provenance, and review status. Preliminary findings include correspondences concerning four floors, timer expiration, and camera exposure threshold; a literal divergence between returning to an exact point and accepting spatial tolerance; and insufficient specification to compare mission duration and hacking-terminal effects numerically. The classifications remain preliminary: the game has not been newly executed, the entire GDD has not been coded, and human review is pending. The contribution is an auditable matrix distinguishing explicit decisions, operational concretization, and documentary gaps. The study neither measures learning nor attributes individual decisions to students or the agent without original records. The findings require verification of the original document, human category review, and subsequent technical checks.

\end{resumotexto}
\noindent\textbf{Keywords:} game design; generative artificial intelligence; traceability; documentary analysis; game development.
\end{minipage}
